\subsection{Khử nhiễu và thứ tự xử lý}
\label{sec:processing-extension}
Nhóm khảo sát khử nhiễu để xem các bộ lọc không gian ảnh hưởng như thế nào đến nhận dạng ảnh thiếu sáng. Nhóm cũng kiểm tra việc đặt bộ lọc trước hoặc sau phép tăng sáng. Các cấu hình được chọn trên validation, còn kết quả trên tập con và toàn bộ test được trình bày ở phần thực nghiệm.

\subsubsection{Câu hỏi và phạm vi khảo sát}
Nhóm đặt ra ba câu hỏi cho khảo sát này. Trước hết, nếu chỉ khử nhiễu mà không tăng sáng, mô hình có nhận dạng tốt hơn không? Tiếp theo, các bộ lọc có cơ chế giữ biên hoặc so sánh mảnh ảnh có giữ được chi tiết phân biệt lớp tốt hơn Gaussian không? Cuối cùng, nhóm muốn biết nên khử nhiễu trước hay sau tăng sáng trong các cấu hình đã chọn. Phạm vi vẫn là ảnh biển báo đã cắt ROI với suy giảm tổng hợp. Nhóm chưa dùng khảo sát này để kết luận cho ảnh toàn cảnh hoặc hệ thống cảnh báo trong xe.

\subsubsection{Cơ chế các bộ lọc và giả thuyết}
Với ảnh nhiễu $I$, bộ lọc Gaussian cho ảnh đầu ra
\begin{equation}
 D_G(I)(p)=\sum_{q\in\mathcal N(p)}G_{\sigma}(p-q)I(q),
 \qquad \sum_{q\in\mathcal N(p)}G_{\sigma}(p-q)=1.
\end{equation}
Trung bình có trọng số làm giảm dao động cục bộ nhưng có thể làm mờ chữ số và đường viền nhỏ trên crop $64\times64$. Bilateral thêm trọng số theo khác biệt màu~\cite{bilateral}:
\begin{equation}
 D_B(I)(p)=\frac{1}{W_p}\sum_{q\in\mathcal N(p)}
 e^{-\|p-q\|^2/(2\sigma_s^2)}
 e^{-\|I(p)-I(q)\|^2/(2\sigma_c^2)}I(q),
\end{equation}
Trong công thức bilateral, $W_p$ là hệ số đưa tổng trọng số về 1. Bộ lọc hạn chế trộn pixel qua các biên có khác biệt màu lớn, nhưng nhiễu mạnh có thể làm sai độ tương tự mà bộ lọc ước lượng. Non-local Means (NLM)~\cite{nlm} dùng dạng $D_N(I)(p)=\sum_q w(p,q)I(q)$. Các trọng số của NLM dựa trên độ giống nhau giữa các mảnh ảnh, thay vì chỉ dựa vào khoảng cách pixel. Khi tăng độ mạnh lọc, phương pháp có thể giảm nhiễu nhiều hơn nhưng cũng xóa chi tiết. Nhóm dùng các hàm OpenCV cho phần triển khai và không huấn luyện thêm mô hình khử nhiễu~\cite{opencv_filter,opencv_nlm}.

Gaussian và bilateral nhận ảnh RGB float32 trong $[0,1]$. Riêng hàm NLM màu của OpenCV yêu cầu ảnh uint8. Vì vậy, mã làm tròn về 8 bit, đổi RGB sang BGR, gọi hàm lọc rồi chuyển kết quả về RGB float32. Hàm còn xử lý màu nội bộ dựa trên không gian Lab. Khi diễn giải PSNR và SSIM, nhóm cần tính đến các bước chuyển đổi này, vì thay đổi đầu ra không chỉ do nguyên lý NLM.

Gọi $E$ là phép tăng sáng và $D$ là phép khử nhiễu. Nhóm so sánh hai chuỗi $E(D(I))$ và $D(E(I))$. Hai chuỗi thường không cho kết quả giống nhau vì phép tăng sáng là phi tuyến và một số bộ lọc phụ thuộc cường độ. Khử nhiễu trước có thể giảm phần nhiễu bị nâng lên khi tăng sáng. Ngược lại, tăng sáng trước có thể làm cấu trúc rõ hơn nhưng cũng khuếch đại sai lệch. Nhóm kiểm tra cả hai thay vì giả định một thứ tự tốt hơn từ đầu.

\subsubsection{Thiết kế chọn tham số công bằng}
Trong khảo sát này, nhóm giữ nguyên cách chia tập theo track, ảnh crop, ba mức suy giảm, noise seed 101/202/303 và ba checkpoint seed 11/22/33. Nhóm không tinh chỉnh CNN bằng ảnh tối hay ảnh đã lọc. Việc chọn tham số vẫn dùng đúng 1.290 ảnh validation của khảo sát tăng sáng. Nhãn test không được dùng để chọn phương pháp. Các cấu hình được lưu trong \path{configs/processing_extension.json}.

\par
\begin{table}[!htbp]
\centering\small
\caption{Lưới tham số khử nhiễu, lựa chọn bằng validation.}
\begin{tabular}{P{3.0cm}P{10.5cm}}
\toprule
Thành phần & Cấu hình và vai trò \\
\midrule
Gaussian & $(k,\sigma)=(3,0{,}5)$ hoặc $(5,0{,}8)$; đối chứng làm trơn cục bộ \\
Bilateral & $d=5$; $(\sigma_c,\sigma_s)=(0{,}05,1{,}5)$ hoặc $(0{,}10,3)$, màu trên thang $[0,1]$ \\
NLM màu & $h=h_{color}=3$ hoặc $7$ trên thang 8 bit; cửa sổ mẫu 3, tìm kiếm 7 \\
Tăng sáng kết hợp & Gamma và CLAHE, giữ tham số đã chọn từ validation đối chứng \\
Thứ tự & Cả khử nhiễu $\to$ tăng sáng và tăng sáng $\to$ khử nhiễu \\
\bottomrule
\end{tabular}
\end{table}\par


Đầu tiên, nhóm chọn một độ mạnh cho mỗi bộ lọc dựa trên Macro-F1 khi khử nhiễu riêng lẻ. Nhóm giữ tham số đó cho cả hai thứ tự kết hợp với từng phép tăng sáng. Có 6 cấu hình riêng lẻ và 12 chuỗi cần đánh giá thêm trên validation. Sau lựa chọn độ mạnh, tập ứng viên test còn 3 bộ lọc riêng lẻ và 12 chuỗi. Cách làm này giúp phân biệt ảnh hưởng của thứ tự với độ mạnh lọc. Tuy nhiên, tham số tốt khi lọc riêng chưa chắc tốt khi kết hợp, vì tăng sáng có thể đổi nhiễu và độ tương tự giữa các vùng. Do đó, kết quả chỉ so sánh các chuỗi với tham số đã cố định, chưa phải kết quả tối ưu toàn cục. Nếu muốn tìm chuỗi tối ưu hơn, nhóm cần chọn đồng thời bộ lọc, bộ tăng cường và thứ tự trên validation, rồi kiểm tra bằng test độc lập. Nhóm giữ Identity, Gamma, CLAHE và quy trình đối chứng đã chọn để so sánh trên cùng ảnh. Nếu có cấu hình trùng nhau, nhóm chỉ tính một lần.

\subsubsection{Khóa subset, full-test và cách phân tích}
Nhóm dùng tập con 430 ảnh để đánh giá toàn bộ họ phương pháp sau khi chọn độ mạnh. Danh sách phương pháp chạy toàn bộ test được xác định \textit{trước khi đọc kết quả tập con}. Trên validation, nhóm chọn cặp bộ lọc và tăng sáng có điểm trung bình hai thứ tự cao nhất, rồi giữ cả hai thứ tự và hai phép riêng lẻ tương ứng. Nhóm bổ sung bộ lọc riêng lẻ tốt nhất, ứng viên tổng thể tốt nhất và các đối chứng. Vì vậy, phần đánh giá toàn bộ test chỉ có tập ứng viên đã sàng lọc, không bao gồm mọi tổ hợp. Hai giao thức được cố định cùng lúc và ghi mã ảnh, tham số, hash mã nguồn cùng hash checkpoint.

Nhóm đánh giá Accuracy, Macro-F1, Weighted-F1, PSNR, SSIM và thời gian xử lý. Với các kết quả nhận dạng, nhóm lấy trung bình noise seed trong từng checkpoint rồi tính trung bình và độ lệch chuẩn qua ba checkpoint. Nhóm còn phân tích F1 từng lớp, ma trận nhầm lẫn và số ảnh đổi từ sai sang đúng hoặc từ đúng sang sai so với Identity. Chênh lệch thứ tự được tính trên cùng ảnh. Thời gian đo chỉ dành cho khối xử lý sau bước chạy làm nóng, chưa phải độ trễ của cả hệ thống camera. Nếu chất lượng ảnh tăng nhưng Macro-F1 giảm, nhóm cần xem lại nét số, đường viền và màu để hiểu kết quả.

Các thuật toán không nhận nhãn hoặc ảnh sáng tham chiếu. Bộ đánh giá mới dùng những thông tin đó để tính điểm. Mọi phương pháp nhận cùng ảnh suy giảm và được kiểm tra bằng cùng bộ phân loại. Khi phân tích thứ tự, nhóm sử dụng các cặp có cùng tham số để tránh nhầm với thay đổi độ mạnh lọc. Các chuỗi mới được đo trong cùng lần chạy và cùng phạm vi tác vụ. Thời gian của đối chứng dùng lại từ lần chạy khác chỉ được trình bày để tham khảo.
